% ****** Start of file template-TIF360-FYM360-blindtext.tex ******
%
% use on Overleaf!!!!
%
\documentclass[%
 reprint,
 amsmath,amssymb,
 aps,
]{revtex4-2}

\usepackage{graphicx}% Include figure files
\usepackage{dcolumn}% Align table columns on decimal point
\usepackage{bm}% bold math
\usepackage{hyperref}% add hypertext capabilities
%\usepackage[mathlines]{lineno}% Enable numbering of text and display math
%\linenumbers\relax % Commence numbering lines
\usepackage{xcolor}
\usepackage{lipsum}


\begin{document}

\title{[C] Context-Aware and Retrieval-Augmented Sarcasm Detection in Social Media}% Force line breaks with \\

\author{Mina Tahmasebi Berjouei}

\date{\today}% It is always \today, today,
             %  but any date may be explicitly specified

\begin{abstract} %%% DO NOT CHANGE!
Sarcasm detection is a challenging natural language processing task because the intended meaning of a text often differs from its literal wording and depends on conversational context. This project investigates whether context-aware modelling can improve sarcasm detection in social media posts. The research question is: \emph{Can a transformer-based classifier that uses both a target comment and its conversational context detect sarcasm more effectively than context-free baselines, and can retrieval-augmented examples or LLM-based reasoning improve interpretability?} The project will use a Reddit sarcasm dataset containing labelled comments and parent-comment context. The main approach will be a supervised transformer classifier, such as RoBERTa or DeBERTa, trained on pairs of parent comments and replies. As comparison methods, I will evaluate simpler machine-learning baselines and an LLM-based zero-shot or few-shot classifier. In addition, I will explore a retrieval-augmented component that retrieves semantically similar sarcastic and non-sarcastic examples to provide extra evidence for prediction and explanation. The expected outcome is a systematic comparison of context-free, context-aware, LLM-based, and retrieval-augmented approaches for sarcasm detection.

\begin{description} %%% DO NOT CHANGE!
\item[Project Topic] %%% DO NOT CHANGE!
{Sarcasm Detection in Social Media.}
\item[Teaching Assistant] %%% DO NOT CHANGE!
{Antonio Ciarlo.}
\end{description} %%% DO NOT CHANGE!
\end{abstract}

\maketitle




\section{\label{sec:intro}Introduction} %%% DO NOT CHANGE!

Sarcasm is a form of figurative language in which the intended meaning of an utterance differs from its literal surface meaning. This makes sarcasm detection difficult for standard sentiment-analysis systems: a sentence may contain positive words while expressing a negative or mocking attitude. In social media, the problem becomes even harder because posts are informal, short, and often depend on previous comments, shared assumptions, and topic-specific background knowledge.

This project studies sarcasm detection as a context-aware NLP classification problem. Instead of classifying only an isolated comment, the project investigates whether using the parent comment and, optionally, retrieved similar examples can improve prediction and explanation. This direction is motivated by prior sarcasm research showing that conversational context is important for both humans and machine-learning systems \cite{khodak2018large}. Therefore, the central research question is whether context-aware transformer models and retrieval-augmented evidence can improve sarcasm detection compared with context-free baselines.


\section{\label{sec:overview}Overview} %%% DO NOT CHANGE!

Sarcasm detection in social media can be formulated as a supervised binary classification problem: given a target comment, predict whether it is sarcastic or non-sarcastic. However, sarcasm is often not identifiable from the target text alone. It may depend on the previous comment, the discussion topic, the speaker's attitude, or an incongruity between the literal sentiment of the reply and the situation described in the conversation. For this reason, this project focuses on context-aware sarcasm detection rather than only context-free text classification.

The dataset considered for the project is based on Reddit comments, where sarcastic comments are self-annotated using the ``/s'' marker. The SARC dataset is especially relevant because it provides not only sarcasm labels but also conversation context, subreddit information, author information, and parent comments \cite{khodak2018large}. This makes it suitable for comparing models that use only the reply with models that also use the parent comment or retrieved contextual examples.

Table~\ref{tab:methodsoverview} summarizes the main approaches that are relevant to the research question. The methods range from simple interpretable baselines to stronger transformer-based models and LLM-based approaches. The selected main method for this project is a supervised context-aware transformer classifier, using the parent comment and target reply as input. Retrieval-augmented and LLM-based methods will be used as additional comparison or explanation components, not as the only classification method.

\begin{table*}
    \caption{{\bf Overview of candidate methods for context-aware sarcasm detection.} The table compares methods that can be used for detecting sarcasm in social media, with emphasis on their relevance to parent-comment context, LLM reasoning, and retrieval-augmented evidence.}
    \label{tab:methodsoverview}
    \begin{tabular}{| l | l | l | l |}
    \hline
       \parbox[t]{2.2cm}{\bf Method} &
       \parbox[t]{3.2cm}{\bf Use case scenario} &
       \parbox[t]{5.7cm}{\bf Advantages and disadvantages} &
       \parbox[t]{5.7cm}{\bf Suitable for this project?} \\
    \hline
        \parbox{2.2cm}{\begin{flushleft}Bag-of-Words / TF-IDF + Logistic Regression\end{flushleft}} &
        \parbox{3.2cm}{\begin{flushleft}Simple context-free sarcasm baseline using only the target comment.\end{flushleft}} &
        \parbox{5.7cm}{\begin{flushleft}Fast, interpretable, and useful as a baseline. However, it cannot model deep semantic meaning, conversation context, or subtle incongruity.\end{flushleft}} &
        \parbox{5.7cm}{\begin{flushleft}Yes, but only as a baseline. It helps measure whether more advanced models actually improve performance.\end{flushleft}} \\
    \hline
        \parbox{2.2cm}{\begin{flushleft}Context-free Transformer\end{flushleft}} &
        \parbox{3.2cm}{\begin{flushleft}Fine-tuning BERT, RoBERTa, or DeBERTa on the target reply only.\end{flushleft}} &
        \parbox{5.7cm}{\begin{flushleft}Captures rich semantic and syntactic information better than n-gram models. The limitation is that it may miss sarcasm that requires the parent comment or topic context.\end{flushleft}} &
        \parbox{5.7cm}{\begin{flushleft}Yes, as a strong neural baseline. It shows how much can be learned from the target text alone.\end{flushleft}} \\
    \hline
        \parbox{2.2cm}{\begin{flushleft}Context-aware Transformer\end{flushleft}} &
        \parbox{3.2cm}{\begin{flushleft}Fine-tuning a transformer on a paired input: parent comment + target reply.\end{flushleft}} &
        \parbox{5.7cm}{\begin{flushleft}Models conversational context and can learn incongruity between what was said before and the reply. It is more computationally expensive and requires careful input formatting.\end{flushleft}} &
        \parbox{5.7cm}{\begin{flushleft}Yes. This is the main selected approach because the research question is specifically about context-aware sarcasm detection.\end{flushleft}} \\
    \hline
        \parbox{2.2cm}{\begin{flushleft}LLM zero-shot / few-shot classifier\end{flushleft}} &
        \parbox{3.2cm}{\begin{flushleft}Prompting an LLM to classify whether a reply is sarcastic, with or without a few examples.\end{flushleft}} &
        \parbox{5.7cm}{\begin{flushleft}Requires little or no task-specific training and can produce natural-language explanations. Disadvantages include cost, latency, possible inconsistency, and weaker reproducibility compared with fine-tuned models.\end{flushleft}} &
        \parbox{5.7cm}{\begin{flushleft}Partly suitable. It is useful for comparison and explanation, but it should not be the only method because the project requires systematic training and evaluation.\end{flushleft}} \\
    \hline
        \parbox{2.2cm}{\begin{flushleft}Retrieval-Augmented Generation (RAG)\end{flushleft}} &
        \parbox{3.2cm}{\begin{flushleft}Retrieving similar sarcastic and non-sarcastic examples or related context before classification.\end{flushleft}} &
        \parbox{5.7cm}{\begin{flushleft}Can provide additional evidence and make predictions more explainable through similar examples. However, RAG is not a classifier by itself and depends strongly on retrieval quality.\end{flushleft}} &
        \parbox{5.7cm}{\begin{flushleft}Suitable as an additional component. It can support an LLM or classifier with retrieved examples, but the main model should remain supervised.\end{flushleft}} \\
    \hline
    \end{tabular}
\end{table*}

\textbf{Bag-of-Words and TF-IDF baselines.} Classical text classification methods represent each comment using word or n-gram frequencies and train a linear classifier such as logistic regression. These methods are useful because they are fast, transparent, and easy to reproduce. They also provide a lower-bound comparison for more advanced models. Their main weakness is that they treat text mostly as surface-level tokens and do not understand context or pragmatic meaning.

\textbf{Transformer-based models.} Pre-trained language models such as BERT, RoBERTa, and DeBERTa are strong candidates for sarcasm detection because they learn contextual word representations from large corpora \cite{devlin2019bert,liu2019roberta}. A context-free transformer can classify the target reply alone, while a context-aware transformer can receive both the parent comment and the reply as one paired input. For example, the input can be structured as: \texttt{[CLS] parent comment [SEP] reply [SEP]}. This is more appropriate for sarcasm because the sarcastic meaning often emerges from the contrast between the reply and the preceding context.

\textbf{LLM-based classification.} Large language models can be used in a zero-shot or few-shot setting by prompting the model to decide whether a comment is sarcastic. This is useful for comparison because LLMs can reason over context and generate explanations. However, LLM-only classification may be expensive, less reproducible, and difficult to evaluate fairly if the prompts are not fixed. Therefore, in this project, the LLM approach will be treated as a comparison method and as an explanation tool rather than the main model.

\textbf{Retrieval-augmented approach.} RAG is usually used for question answering, but it can be adapted to this project by retrieving semantically similar training examples for a new target example \cite{lewis2020retrieval}. For instance, the system can embed the parent-comment and reply pair, retrieve similar sarcastic and non-sarcastic examples, and provide these examples to an LLM or use them as additional evidence. This is relevant because sarcasm is often pattern-based and context-dependent. Nevertheless, RAG should not replace supervised classification; it should be used to improve interpretability or to test whether retrieved examples help the prediction.

\textbf{Selected approach.} The main method selected for the project is a supervised context-aware transformer classifier. The model will be trained to classify a reply as sarcastic or non-sarcastic using both the reply and its parent comment. The project will compare this model against a simple TF-IDF baseline, a context-free transformer baseline, and an LLM-based zero-shot or few-shot baseline. If time permits, a retrieval-augmented variant will be added to retrieve similar examples and produce explanations for selected predictions. This design is appropriate because it directly addresses the central research question: whether conversational context and retrieval-based evidence improve sarcasm detection beyond target-text-only classification.














\begin{thebibliography}{9}

\bibitem{khodak2018large}
M. Khodak, N. Saunshi, and K. Vodrahalli, ``A Large Self-Annotated Corpus for Sarcasm,'' in \emph{Proceedings of the Eleventh International Conference on Language Resources and Evaluation}, 2018.

\bibitem{devlin2019bert}
J. Devlin, M.-W. Chang, K. Lee, and K. Toutanova, ``BERT: Pre-training of Deep Bidirectional Transformers for Language Understanding,'' in \emph{Proceedings of NAACL-HLT}, 2019.

\bibitem{liu2019roberta}
Y. Liu et al., ``RoBERTa: A Robustly Optimized BERT Pretraining Approach,'' \emph{arXiv preprint arXiv:1907.11692}, 2019.

\bibitem{lewis2020retrieval}
P. Lewis et al., ``Retrieval-Augmented Generation for Knowledge-Intensive NLP Tasks,'' in \emph{Advances in Neural Information Processing Systems}, 2020.

\end{thebibliography}


\end{document}
